\chapter{Podstawy teoretyczne}

\section{Architektura systemu Android}
Android jest otwartą, wielowarstwową platformą opartą na jądrze Linuksa. Stos oprogramowania obejmuje m.in.\ jądro (kernel) z mechanizmami bezpieczeństwa (np.\ SELinux), warstwę HAL (Hardware Abstraction Layer), biblioteki natywne, środowisko uruchomieniowe ART oraz framework Java/Kotlin z zestawem API, na którym budowane są aplikacje i usługi systemowe \cite{AOSPArchitecture,AndroidPlatformArchitecture}. W nowszych wersjach platformy pojawiły się dodatkowe izolatory, takie jak \emph{SDK sandbox} (separowany proces dla wybranych SDK), co jeszcze bardziej ogranicza wpływ komponentów zewnętrznych na proces aplikacji \cite{SdkSandbox}.

\section{Model bezpieczeństwa}
Fundamentem bezpieczeństwa Androida jest \emph{piaskownica} (application sandbox): każda aplikacja otrzymuje unikalny UID i działa w odseparowanym procesie z izolowanym przestrzennie systemem plików. Dodatkową ochronę zapewnia polityka MAC (SELinux), a także model uprawnień ograniczający dostęp do wrażliwych zasobów (lokalizacja, kamera, mikrofon) – w tym uprawnienia żądane w czasie działania (\emph{runtime permissions}) \cite{AndroidSandbox,AndroidPermissionsOverview,AndroidPermissionRequest}. Łańcuch zaufania domykają mechanizmy podpisywania aplikacji, Verified Boot/AVB oraz proces weryfikacji w ekosystemie sklepu (Play Protect, audyty polityk), co Google cyklicznie raportuje \cite{GooglePlaySafety2024}. 
W kontekście komunikacji sieciowej część aplikacji implementuje \emph{certificate pinning}, które utrudnia badaczowi pasywne przechwycenie treści HTTPS przez klasyczne proxy – poprawnie wdrożony pinning ma powodować odrzucenie połączenia z certyfikatem wystawionym przez zaufane, lecz nieprawidłowe (dla danej domeny) CA \cite{OWASPMASTGPinning}.

\section{Protokoły}
Przeglądane aplikacje korzystają z rodziny protokołów HTTP w różnych wersjach oraz z TLS do uwierzytelnienia i poufności. 
\begin{itemize}
  \item \textbf{HTTP/1.1} – klasyczny, tekstowy protokół nad TCP (przykładowe różnice semantyczne i transportowe względem nowszych wersji omawiają poniższe punkty HTTP/2/HTTP/3). 
  \item \textbf{HTTP/2} – binarne ramkowanie, kompresja nagłówków i \emph{multiplexing} wielu strumieni na jednym połączeniu TCP w celu redukcji opóźnień \cite{RFC9113}.
  \item \textbf{HTTP/3} – mapowanie semantyki HTTP na QUIC (UDP), eliminujące wpływ HoL na poziomie transportu i skracające zestawienie połączenia \cite{RFC9114,RFC9000}. 
  \item \textbf{TLS 1.3} – domyślny protokół kryptograficzny dla HTTPS, skracający \emph{handshake}, upraszczający zestaw szyfrów i wspierający 0-RTT (z zastrzeżeniami bezpieczeństwa) \cite{RFC8446}.
\end{itemize}

\section{Techniki analizy ruchu}
W analizach bezpieczeństwa mobilnego wykorzystuje się zwykle dwa komplementarne podejścia: 
\emph{(1) przechwytywanie i inspekcję ruchu} oraz \emph{(2) instrumentację/dynamiczną obserwację aplikacji}. 
Do przechwytywania pakietów i strumieni wykorzystywane są m.in.\ \textit{tcpdump} i \textit{Wireshark}, a do pośredniczenia w ruchu HTTP(S) – \textit{mitmproxy} (instalacja zaufanego CA na urządzeniu w celu deszyfracji ruchu) \cite{tcpdump,WiresharkDocs,MitmproxyDocs,OWASPMASTGIntercept}. 
W przypadku aplikacji z poprawnym \emph{certificate pinning} wymagane bywa podejście z instrumentacją (np.\ \textit{Frida}) lub środowisko z większą kontrolą (urządzenie z dostępem \textit{root}/emulator, własny trust store), przy czym nowoczesne techniki (np.\ niestandardowi \emph{trust managerzy}, natywne weryfikacje, integracja z \emph{Play Integrity}) istotnie podnoszą poprzeczkę dla badań dynamicznych \cite{Frida,OWASPMASTGPinning}. 
Wybór między emulatorem a urządzeniem fizycznym wpływa na wierność zachowania (emulator ułatwia kontrolę i automatyzację; urządzenie fizyczne lepiej oddaje ograniczenia sprzętowe i integracje OEM) \cite{OWASPMASTGIntercept}.

\section{Prywatność i tracking}
Problem prywatności w aplikacjach mobilnych obejmuje zarówno zakres i uzasadnienie uprawnień, jak i obecność komponentów stron trzecich (SDK analityczne/reklamowe). Do statycznego wykrywania trackerów i przeglądu uprawnień w APK wykorzystywane są m.in.\ \textit{Exodus Privacy} oraz narzędzia deweloperskie Androida \cite{ExodusPrivacy,AndroidPermissionsOverview}. Równolegle platforma rozwija rozwiązania mające ograniczyć śledzenie międzyaplikacyjne (np.\ \emph{Privacy Sandbox on Android}) \cite{PrivacySandboxAndroid}. W praktyce ocena prywatności powinna łączyć analizę statyczną i dynamiczną (ruch sieciowy, zachowanie w tle), a także kontrolować zużycie zasobów (bateria/CPU) – symptomy nadmiernej telemetrii mogą ujawniać się jako wzrost aktywności sieciowej w tle czy anomalia energetyczne.
